\subsubsection{Brandy Flip}
\noindent
\begin{minipage}{.49\textwidth}
    \noindent\textbf{\LARGE Brandy Flip\index{Brandy Flip}}
\end{minipage}
\hspace{.02\textwidth}
\begin{minipage}{.49\textwidth}
    \raggedleft
    \textbf{\large Brandy}\\
    \sweettartguage{-0.3}\index[drinks]{Sweet!30\%!Brandy Flip}
\end{minipage}
\vspace{.5em}\hrule\vspace{.5em}
\noindent
\textbf{\large Flavor Profile:} Sweet\index[flavor-profile]{Sweet}, Rich\index[flavor-profile]{Rich}, Thick\index[flavor-profile]{Thick}, Reminds you of Christmas\index[flavor-profile]{Christmas}.
\vspace{.5em}\hrule\vspace{.5em}
\noindent
\textbf{\large History:} The term "flip" dates back to the late 1600s and refers to a drink consisting of beer, rum, eggs, and sugar. Flipping occurred by adding a red-hot iron to the drink causing it to froth or "flip". The version known today was first seen in Jerry Thomas's book called the Bartender's Guide. It is typically made with brandy or sherry. The cocktail is also very closely related to eggnog, simply add cream or milk!
\vspace{.5em}\hrule\vspace{.5em}
\noindent
\begin{minipage}[t]{.49\textwidth}
    \noindent
    \textbf{\large Ingredients}
    \begin{itemize}
        \item Cognac 60ml\index[ingredient]{Cognac}
        \item Demerara Sugar 22,5ml\index[ingredient]{Demerara Sugar}
        \item Egg 1\footnote{Consuming raw eggs may increase risk of foodborne illness.}\index[ingredient]{Egg}
    \end{itemize}
    \noindent
    \textbf{\large Garnish}
    \begin{itemize}
        \item Ground Nutmeg\index[ingredient]{Nutmeg}
    \end{itemize}
\end{minipage}
\hspace{.01\textwidth}\vrule\hspace{.01\textwidth}
\begin{minipage}[t]{.49\textwidth}
    \noindent
    \textbf{\large Tools/Glassware}
    \begin{itemize}
        \item Shaker
        \item Strainer \& Fine Strainer (optional)
        \item Jigger
        \item Small stemmed glass\index[glass]{Small stemmed glass} or 120ml coupe\index[glass]{Coupe}
        \item Microplane
    \end{itemize}
\end{minipage}
\vspace{.5em}\hrule\vspace{.5em}
\noindent
\begin{minipage}{\textwidth}
    \noindent
    \textbf{\large Instructions}
    \begin{itemize}
        \item Place all ingredients into a shaker.
        \item Dry shake for 20 seconds (no ice).
        \item Add ice and shake for 8--10 seconds.
        \item Double strain into small stemmed glass or coupe.
        \item Garnish with fresh ground nutmeg.
    \end{itemize}
\end{minipage}
\vspace{.5em}\hrule\vspace{.5em}
\noindent
\textbf{\large Fun Tip:} To make eggnog use the same recipe, and add 30ml of heavy cream or 60ml of milk!
\newpage
